% from ega2-1-fr.tex
\oldpage[II]{113}

pour espace sous-jacent $Z$ (\hyperref[I.5.2.1-fr]{I, 5.2.1}), et il résulte
donc de (\hyperref[II.6.1.5-fr]{6.1.5, (i) et (ii)}) qu'on peut se borner à
prouver que $f(X)$ est fermé dans $Y$. En vertu de
(\hyperref[II.6.1.5-fr]{6.1.5, (vii)}), on peut supposer $X$ et $Y$ réduits~;
en outre, si $T$ est le sous-préschéma fermé réduit de $Y$ ayant pour espace
sous-jacent $\overline{f(X)}$ (\hyperref[I.5.2.1-fr]{I, 5.2.1}), on sait que
$f$ se factorise en $X\xrightarrow{g}T\xrightarrow{j}Y$, où $j$ est le
morphisme d'injection (\hyperref[I.5.2.2-fr]{I, 5.2.2}), et comme $j$ est
séparé (\hyperref[I.5.5.1-fr]{I, 5.5.1, (i)}), il résulte de
(\hyperref[II.6.1.5-fr]{6.1.5, (v)}) que $g$ est un morphisme entier. On peut
donc supposer que $f(X)$ est dense dans $Y$. Enfin, la question étant locale
sur $Y$, on peut se borner au cas où $Y=\operatorname{Spec}(A)$. Alors
$X=\operatorname{Spec}(B)$, où $B$ est une $A$-algèbre entière sur $A$
(\hyperref[II.6.1.4-fr]{6.1.4})~; en outre $A$ est réduit
(\hyperref[I.5.1.4-fr]{I, 5.1.4}), et l'hypothèse que $f(X)$ est dense dans
$Y$ entraîne que l'homomorphisme $\varphi:A\to B$ correspondant à $f$ est
injectif (\hyperref[I.1.2.7-fr]{I, 1.2.7}). Dire que dans ces conditions
$f(X)=Y$ signifie que tout idéal premier de $A$ est la trace sur $A$ d'un
idéal premier de $B$, ce qui n'est autre que le premier théorème de
Cohen-Seidenberg (\cite[t.~I, p.~257, th.~3]{I-13}).

\begin{corollary}[6.1.11]
\label{II.6.1.11-fr}
Tout morphisme fini $f:X\to Y$ est projectif.
\end{corollary}

En effet, comme $f$ est affine, $\mathcal{O}_X$ est un
$\mathcal{O}_X$-Module très ample relativement à $f$
(\hyperref[II.5.1.2-fr]{5.1.2})~; en outre $f_*(\mathcal{O}_X)$ est un
$\mathcal{O}_Y$-Module quasi-cohérent de type fini
(\hyperref[II.6.1.2-fr]{6.1.2})~; enfin $f$ est séparé, de type fini et
universellement fermé (\hyperref[II.6.1.10-fr]{6.1.10}), et on est donc dans
les conditions d'applications du critère
(\hyperref[II.5.5.4-fr]{5.5.4, (ii)}).

\begin{proposition}[6.1.12]
\label{II.6.1.12-fr}
Soient $f:X'\to X$ un morphisme fini, et soit
$\mathcal{B}=f_*(\mathcal{O}_{X'})$ (qui est une $\mathcal{O}_X$-Algèbre
quasi-cohérente et un $\mathcal{O}_X$-Module de type fini). Soit
$\mathcal{F}'$ un $\mathcal{O}_{X'}$-Module quasi-cohérent~; pour que
$\mathcal{F}'$ soit localement libre, de rang $r$, il faut et il suffit que
$f_*(\mathcal{F}')$ soit un $\mathcal{B}$-Module localement libre de rang
$r$.
\end{proposition}

Il est clair que si $f_*(\mathcal{F}')|U$ est isomorphe à
$\mathcal{B}^r|U$ ($U$ ouvert de $X$), $\mathcal{F}'|f^{-1}(U)$ est isomorphe
à $\mathcal{O}_{X'}^r|f^{-1}(U)$
(\hyperref[II.1.4.2-fr]{1.4.2}). Inversement, supposons que $\mathcal{F}'$
soit localement libre de rang $r$ et montrons que $f_*(\mathcal{F}')$ est
localement isomorphe à $\mathcal{B}^r$ en tant que $\mathcal{B}$-Module. Soit
$x$ un point de $X$~; lorsque $U$ parcourt un système fondamental de
voisinages affines de $x$, $f^{-1}(U)$ parcourt un système fondamental de
voisinages affines (\hyperref[II.1.2.5-fr]{1.2.5}) de l'ensemble fini
$f^{-1}(x)$, puisque $f$ est fermé
(\hyperref[II.6.1.10-fr]{6.1.10}). La proposition résulte alors du lemme
suivant~:

\begin{lemma}[6.1.12.1]
\label{II.6.1.12.1-fr}
Soient $Y$ un préschéma, $\mathcal{E}$ un $\mathcal{O}_Y$-Module localement
libre de rang $r$, $Z$ une partie finie de $Y$ contenue dans un ouvert affine
$V$. Il existe alors un voisinage $U\subset V$ de $Z$ tel que
$\mathcal{E}|U$ soit isomorphe à $\mathcal{O}_Y^r|U$.
\end{lemma}

On peut évidemment supposer $Y$ affine~; pour tout $z_i\in Z$, il existe dans
l'adhérence $\overline{\{z_i\}}$ au moins un point fermé $z'_i$
(\hyperref[0.2.1.3-fr]{0, 2.1.3})~; si $Z'$ est l'ensemble des $z'_i$, tout
voisinage de $Z'$ est un voisinage de $Z$, et on peut donc supposer $Z$ fermé
dans $Y$ et discret. Considérons le sous-préschéma fermé réduit de $Y$ ayant
$Z$ pour espace sous-jacent (\hyperref[I.5.2.1-fr]{I, 5.2.1}) et soit
$j:Z\to Y$ l'injection canonique~; $j^*(\mathcal{E})=\mathcal{E}\otimes_Y
\mathcal{O}_Z$ est localement libre de rang $r$ sur le schéma discret $Z$,
et par suite est isomorphe à $\mathcal{O}_Z^r$~; autrement dit il existe $r$
sections $s_i$ ($1\leq i\leq r$) de $\mathcal{E}\otimes_Y\mathcal{O}_Z$
au-dessus de $Z$ telles que l'homomorphisme
$\mathcal{O}_Z^r\to\mathcal{E}\otimes_Y\mathcal{O}_Z$ défini par ces sections
soit bijectif. Mais $Y=\operatorname{Spec}(A)$ est affine, $Z$ est défini par
un idéal $\mathfrak{J}$ de $A$ et on a $\mathcal{E}=\widetilde{M}$, où $M$
est un $A$-module~; les $s_i$ sont des éléments

